%==============================================================================
\section{The UBench Framework and Protocol} \label{sec:metodologia}
%==============================================================================

%------------------------------------------------------------------------------
\subsection{Framework overview} \label{sec:visao_geral_ubench}

UBench is a modular Python/PyTorch framework executed end-to-end by a single command (Fig.~\ref{fig:pipeline_ubench}).
An orchestration module probes the GPU, runs preprocessing when its versioned output is missing or outdated, builds subject-partitioned data loaders, trains all architectures within \textbf{one shared training loop}, and runs a common benchmark module.
Since no architecture has its own training, metric, or timing path, the resulting numbers are comparable by design.
Each run receives an identifier that isolates its artifacts and a metadata file recording the Git commit hash and a dirty-tree flag, library and CUDA versions, GPU, seed and determinism flag, effective configuration, partitions, and the hash of the dependency lock file.
Interrupted runs resume from the most recent checkpoint, and a run is marked successful only if no training failure is recorded.

\begin{figure}[tbp]
\centering
\begin{tikzpicture}[
  >=latex, thick,
  every node/.style={font=\scriptsize},
  proc/.style={rectangle, rounded corners=2pt, draw, fill=black!12, align=center, text width=2.45cm, minimum height=0.95cm, inner sep=2pt},
  art/.style={rectangle, draw, fill=black!3, align=center, text width=2.45cm, minimum height=0.95cm, inner sep=2pt},
  node distance=0.35cm and 0.32cm]
 \node[art]  (raw)  {Raw 16-bit TIFFs\\+ landmarks + boxes};
 \node[proc, right=of raw] (pre) {Preprocessing: \textdegree{}C, crop, mask rasterization, $256^2$};
 \node[art, right=of pre] (proc) {Frozen corpus\\+ versioned manifest};
 \node[proc, below=of proc] (cv) {\texttt{GroupKFold} by subject ($K{=}5$)};
 \node[proc, left=of cv] (train) {Shared training loop\\(recipe per family)};
 \node[art, left=of train] (models) {U-Net\\TransUNet\\\swinp};
 \node[proc, below=of models] (bench) {Benchmark: accuracy, latency, memory};
 \node[proc, right=of bench] (stats) {Paired statistics by subject};
 \node[art, right=of stats] (out) {Run artifacts\\+ provenance metadata};
 \draw[->] (raw) -- (pre); \draw[->] (pre) -- (proc); \draw[->] (proc) -- (cv);
 \draw[->] (cv) -- (train); \draw[<->] (train) -- (models);
 \draw[->] (models) -- (bench); \draw[->] (bench) -- (stats); \draw[->] (stats) -- (out);
\end{tikzpicture}
\caption{UBench pipeline. Rectangles: persisted artifacts; rounded boxes: processing stages.}
\label{fig:pipeline_ubench}
\end{figure}

%------------------------------------------------------------------------------
\subsection{Dataset and reference masks} \label{subsec:charlotte_thermalface}

Charlotte-ThermalFace \cite{Ashrafi2022} contains about 10{,}000 thermal face images of 10 subjects, acquired at ambient temperatures from $20.5$ to $26.5$\,\degC{} with varying humidity, air velocity, and capture distance, and stored as 16-bit raw radiometric frames.
After discarding unreadable frames, missing landmarks, and degenerate crops, the corpus has \textbf{8{,}074 images}.
The dataset provides anatomical \textbf{landmarks}, not segmentation masks; reference masks are therefore \textbf{derived}: each image is cropped around the face bounding box with a fixed margin, landmarks are shifted to the (border-clamped) crop origin, and the landmarks of each region are grouped into a closed polygon and rasterized, following the precedent of Joshi et al.~\cite{Joshi2022}.
The 10 classes are background, lower face contour, left and right eyebrows, nose, left and right eyes, mouth, lips, and forehead; left/right follow the observer's viewpoint.
The derivation code is deterministic and version-controlled, and the generated masks are frozen and verified by checksum.
Derived masks impose an architecture-independent \textbf{performance ceiling}: small regions bounded by few landmarks (eyebrows, lips) suffer proportionally larger polygonal approximation error, and no manual masks exist to quantify this noise.

%------------------------------------------------------------------------------
\subsection{Radiometric preprocessing and augmentation} \label{sec:preprocessamento_radiometrico}

Raw values (centikelvin) are converted to degrees Celsius and normalized using a \textbf{fixed} range, $[20, 40]$\,\degC{}$\,\mapsto[0,1]$, identical for all images.
The usual per-image min--max normalization would eliminate absolute temperature information, which is exactly what distinguishes thermography from a grayscale image; the fixed range preserves comparability across images and subjects, as advocated by \cite{Joshi2022}.
Images are resized to $256\times256$ with bilinear interpolation and masks with nearest neighbor.
A manifest records the schema version, unit, resolution, and normalization; the loader raises an explicit error if it is missing or incompatible, so data produced under an older convention cannot be silently reinterpreted.
Because pixel values are temperatures, photometric augmentations (brightness, contrast) are excluded.
Augmentation combines geometric transforms (horizontal flip and affine), applied jointly to image and mask, with physically plausible radiometric perturbations applied in \degC{}: an additive sensor offset of up to $\pm0.5$\,\degC{} and Gaussian noise.

%------------------------------------------------------------------------------
\subsection{Evaluated architectures} \label{sec:arquiteturas_avaliadas}

All architectures are trained from scratch on the thermal data only (Table~\ref{tab:arquiteturas}).
\textbf{U-Net} \cite{Ronneberger2015} follows the canonical topology (widths $64\!\rightarrow\!512$, 1024-channel bottleneck).
\textbf{TransUNet} \cite{Chen2021} uses a convolutional encoder, a ViT-B/16-style bottleneck (768-d embeddings, 12 blocks, 12 heads) with a learned absolute positional embedding, and a convolutional decoder from $H/16$ to full resolution.

\textbf{\swinp{}} retains a hierarchical encoder consistent with the Swin-T scale \cite{Liu2021Swin} ($d{=}96$, depths $[2,2,6,2]$, heads $[3,6,12,24]$, $8\times8$ windows, $4\times4$ patches with patch merging) but \textbf{does not reproduce} Swin-UNet++ \cite{Liu2021Materials}, hence its name (abbreviated \swint{} in tables).
Four simplifications were necessitated by the 6~GB GPU memory constraint: a Swin decoder with patch expanding, a full nested skip grid, and deep supervision would exceed the available memory budget at $256\times256$ without mixed precision, so the decoder is convolutional with collapsed nested skips and a $4\times$ interpolation from $H/4$.
Three additional omissions affect the attention mechanism itself: the shifted-window mask, the relative position bias, and any positional encoding.
Without them, window attention is permutation-invariant within each window.
Its results are therefore reported as a \textbf{simplified baseline}, not as evidence about Swin-style window attention.
Pre-trained variants (SwinV2-T and R50+ViT-B/16 from \texttt{timm}) exist in the pipeline but are excluded by design: mixing ImageNet-pretrained and from-scratch models would primarily reflect the effect of pre-training.

\begin{table}[tbp]
\centering
\caption{Evaluated architectures (none pre-trained)}
\label{tab:arquiteturas}
\footnotesize
\setlength{\tabcolsep}{3pt}
\begin{tabular}{@{}>{\raggedright\arraybackslash}p{1.55cm}>{\raggedright\arraybackslash}p{2.75cm}>{\raggedright\arraybackslash}p{2.35cm}r@{}}
\toprule
\textbf{Architecture} & \textbf{Encoder} & \textbf{Attention} & \textbf{Par.\,(M)} \\
\midrule
U-Net & Convolutional & None & 31.0 \\
TransUNet & Conv.\ + ViT-B/16 bottleneck & Global, at bottleneck & 105.7 \\
\swint{} & Hierarchical, Swin-T scale & Windowed, simplified & 40.9 \\
\bottomrule
\end{tabular}
\end{table}

%------------------------------------------------------------------------------
\subsection{Training protocol} \label{subsec:recipes_otimizacao}

\textbf{Subject-disjoint cross-validation.}
We use $K=5$ folds with \texttt{GroupKFold} grouped by subject identity; an automated test asserts that no subject appears in both training and validation of any fold.
Each subject is validated exactly once.
\texttt{GroupKFold} balances sample counts and is deterministic, yielding the validation pairs \{S2,S5\}, \{S8,S9\}, \{S6,S10\}, \{S3,S7\}, and \{S1,S4\} for folds 1--5.

\textbf{Family-specific recipes and budget.}
Imposing one optimizer and schedule on all families is not neutral; it is the poorly tuned baseline that \cite{Isensee2024} identify as a source of spurious conclusions.
Each family therefore uses its canonical recipe, without any hyperparameter search (Table~\ref{tab:hiperparametros}): the transformer family uses AdamW with decoupled weight decay \cite{Loshchilov2019}, linear warm-up, and cosine decay, as in \cite{Liu2021Swin}; U-Net uses Adam with \texttt{ReduceLROnPlateau}, the conventional setting for medical U-Nets \cite{Siddique2021}.
The budget is \textbf{100 epochs} for every architecture.
Batch sizes are selected by the hardware probe as the largest that fits in memory (8 for U-Net, 6 for the transformers), so the transformers perform about $33\%$ more weight updates per epoch (1{,}077 vs.\ 807); this asymmetry, like gradient clipping, favors the transformer family and is declared rather than hidden.
All models minimize $\mathcal{L}=0.5\,\mathcal{L}_{CE}+0.5\,\mathcal{L}_{Dice}$ \cite{Milletari2016,Sudre2017,Jadon2020} with uniform class weights.

\begin{table}[tbp]
\centering
\caption{Training hyperparameters by family}
\label{tab:hiperparametros}
\footnotesize
\setlength{\tabcolsep}{3pt}
\begin{tabular}{@{}lll@{}}
\toprule
 & \textbf{U-Net (CNN)} & \textbf{TransUNet, \swint{}} \\
\midrule
Optimizer & Adam & AdamW \\
Learning rate & $1{\times}10^{-4}$ & $1{\times}10^{-4}$ \\
Weight decay & $0$ & $0.05$ \\
Scheduler & Plateau (pat.\ 5, $\times0.5$) & 5\% warm-up $\rightarrow$ cosine \\
Grad.\ clipping & none & $1.0$ \\
Batch / epochs & 8 / 100 & 6 / 100 \\
\bottomrule
\end{tabular}
\end{table}

\textbf{Epoch selection.}
The \textbf{primary} result uses the model at the \textbf{final epoch} of the budget, fixed \textit{a priori} and independent of validation performance.
The \textbf{secondary} result uses the epoch with the best validation mIoU, which is optimistic by construction because the same partition selects and reports.
The plateau scheduler of U-Net is itself conditioned on validation mIoU, a mild adaptation that favors U-Net, i.e., acts against the main result.

%------------------------------------------------------------------------------
\subsection{Metrics and cost instrumentation} \label{subsec:metricas_acuracia}

\textbf{Accuracy.} mIoU and Dice are computed on argmax predictions from a single confusion-matrix implementation (\texttt{torchmetrics}) shared by the training loop and the benchmark module, which rules out silent discrepancies between two implementations of the same metric.
The per-image macro average excludes classes absent from the reference, and the background is reported separately.
Boundary metrics (HD95, NSD) are out of scope.

\textbf{Cost.} Latency discards 20 warm-up iterations, times 200, and explicitly synchronizes the GPU before reading the clock; otherwise one measures the asynchronous kernel queue rather than the computation.
Latency and peak memory use a \textbf{fixed batch of 4} for all architectures, since peak memory scales with batch size.
Only one model resides on the GPU, and \texttt{reset\_peak\_memory\_stats} and \texttt{max\_memory\_allocated} bracket the forward pass.
Three repetitions give identical values, and a fourth with 300~MB of unrelated tensors resident differs by less than 1~MB, so the measurement is independent of evaluation order.
We also report parameters, checkpoint size, and wall-clock training time per fold.

%------------------------------------------------------------------------------
\subsection{Statistical analysis} \label{sec:analise_estatistica}

Comparisons are \textbf{paired by subject}: each subject is evaluated once by a model that never saw it, giving $n=10$ paired observations per comparison.
We report mean differences, wins, the exact two-sided Wilcoxon signed-rank test, 95\% bootstrap confidence intervals of the paired differences (20{,}000 resamples), and Holm correction \cite{Holm1979}.
The pairing unit matters: the smallest attainable exact $p$-value is $2/2^{n}$, i.e., $0.0625$ when pairing by fold ($n=5$), above the 5\% threshold however consistent the folds are, but $\approx0.0020$ for $n=10$.
Because the two subjects of a fold share one trained model, subject-level $p$-values are slightly optimistic, so fold-level pairing ($n=5$ independent models) is reported as a secondary analysis.

%------------------------------------------------------------------------------
\subsection{Execution environment and pipeline audit} \label{subsec:auditoria_defeitos}

The reported run used one NVIDIA GTX~1660~Ti (6~GB), PyTorch~2.5.1 with CUDA~12.1, seed~42 with deterministic mode, and no mixed precision (unsupported on this GPU class), totaling \textbf{186.9 GPU-hours} for the 15 architecture--fold combinations.
The run was launched from a working tree with uncommitted changes outside the pipeline code; the source tree is byte-identical to the tagged commit, and the difference is versioned as a patch, so the run is reproducible from a commit plus an inspectable patch.
Before the run, the pipeline underwent a systematic audit.
Overall, \textbf{32 issues} were cataloged: 26 found by the pre-run audit, four revealed by the pattern of the results (UB-27 to UB-30), and two by verifying the corrections; 28 are corrected and verified, mostly by automated regression tests, and four are disclosed exceptions (e.g., UB-17, the attention omissions of \swinp{}).
The pre-run classes covered divergent metric definitions, partition integrity, preprocessing and architectural fidelity, efficiency measured under unequal conditions, and environment drift.
Issues relevant to interpreting the results are discussed where they arise.
